
il
suffit, d'après la définition de @@M0@@ comme adjoint à gauche du foncteur image directe, de définir un homomorphisme @@M1@@-linéaire
@@M2@@
Pour ce faire, il faut se donner pour tout ouvert @@M3@@ de @@M4@@ un homomorphisme :
@@M5@@
compatible aux restrictions aux ouverts de @@M6@@, et prendre l'homomorphisme correspondant sur les faisceaux associés. L'homomorphisme (2.5.2) est défini de la façon suivante : soit @@M7@@, où, @@M8@@ étant de degré @@M9@@, @@M10@@; on associe alors à @@M11@@ la section @@M12@@ de @@M13@@ sur @@M14@@, en notant @@M15@@ l'homomorphisme canonique @@M16@@ (cf. EGA II 2.6.3 et 2.6.4). On vérifie sans peine les compatibilités nécessaires, d'où par recollement l'homomorphisme (2.5.1).

Il est clair que dans le cas où @@M17@@ est quasi-cohérent l'homomorphisme ainsi construit est celui de EGA II (3.3.4.1).

L'homomorphisme de foncteur @@M18@@ sur la catégorie des @@M19@@-Modules donne par passage aux catégories dérivées un homomorphisme de foncteur sur @@M20@@ :
@@M21@@

\textbf{Proposition 2.6.} \uline{Soient} @@M22@@ \uline{un schéma quasi}-\uline{compact}, @@M23@@ \uline{un} @@M24@@\uline{-Module localement libre de type fini}, @@M25@@, \uline{et} @@M26@@ \uline{un complexe parfait sur} @@M27@@. \uline{Alors} \uline{il existe un entier} @@M28@@ \uline{tel que pour} @@M29@@ \uline{les propriétés suivantes soient vérifiées} :
